\scanpage{35}
\mspage{5}

\uline{L5} Si $X,Y$ sont \inserted{Cat-}\unsure{fibrées|fibrés} sur $S$, et $X\unsure{\text{---}|\to}Y$ est
un $S$-morphisme \uline{cartésien}, induisant sur les
fibres des $X_s\longrightarrow Y_s$ qui sont dans $W$, alors
$f$ est dans $W$.

\deleted{C'est la} \inserted{Pour $\deleted{\unsure{D}}W_{\mathbb D}$,} prop. 6. page 71 de $\mathrm I$.

Ces axiomes $L_1$ à L5 (équivalents à
L1, L3$'$ ($\mathbb{\Delta}_1\to e$ est dans $W$) et L5) sont dans
la nature d'axiomes de stabilité, assurant
que $W$ est assez grand. Il y a une autre
condition qui limite $W$ à ne pas être trop
grand :

L a) Si $f\in W$, alors $\pi_0(f)$ est bijectif.

Ceci n'est pas satisfait pour tous les $W_{\mathbb D}$
– p.ex. si $D$ est le dérivateur \uline{nul}, alors \deleted{$W_{\mathbb D}$}
$W_{\mathbb D}=\mathrm{Fl}\,\mathrm{Cat}$ ! Mais pour les dérivateurs \unsure{utiles|s'utilisant}
en pratique, cela semble toujours satisfait

L b) c) Le foncteur
\[
\mathrm{Cat}\text{---}\mathrm{Cat}.\,W^{-1}
\]
commute aux sommes \inserted{finies} et aux produits finis$^{*}$.

% slanted
\margin{$^{*}$\unsure{Ces|Les} \unsure{axiomes} si $(\mathrm{Cat},W)$ provient d'un triple de Quillen (avec des Q-fibrations, Q-équiv. faibles, Q-cofibrations)}

À examiner ce que ça signifie en termes
concrets sur $W$ – je n'ai pas vérifié que
c'est OK pour les $W_{\mathbb D}$.

L d) Si $f:X\to Y$ dans $W$, alors $f^\circ:X^\circ\longrightarrow Y^\circ$ aussi,
C'est le cor. 2 page 75 de $(\mathrm I)$,

% archive sheet number: 34
